
\subsection{Analysis of \model{}}
% \section{Ablations \YA{Rather Analysis of SteerViT (since replacing with different backbone is not ablation, neither is scaling} }
\label{sec:ablations}

% \begin{table}[t]
% \centering
% \footnotesize
% \setlength{\tabcolsep}{6pt}
% \begin{threeparttable}
% \caption{%
% Architecture ablations on DINOv2 ViT-B/14.
% Each row modifies one component from the full \model{} configuration.
% % \MT{PODS is 58.3 I think .. there may be typos in this table and then corresponding text}
% % \MG{No 58.3 is with MLLM descriptions, 53.4 is for human generated ones}
% }
% \label{tab:arch_ablation}
% \begin{tabular}{@{}l
%                 S[table-format=2.1]
%                 S[table-format=2.1]
%                 S[table-format=2.1]@{}}
% \toprule
% {Variant}
%   & {FG-CLS $\uparrow$}
%   & {Steer.\ $\uparrow$}
%   & {PODS $\uparrow$} \\
% \midrule
% \color{gray} DINOv2 (frozen)
%   & \color{gray} 89.0
%   & \color{gray} 43.5
%   & \color{gray} 28.6 \\
% \rowcolor{\methodcolor}
% \textbf{\model{} (ours)}
%    & \textbf{76.3}
%    & \textbf{95.9}
%    & \textbf{53.4} \\
% \addlinespace
% \rowcolor{gray!20}
% \multicolumn{4}{l}{Architecture ablations} \\
% Late fusion
%   & 86.5
%   & 93.4
%   & 45.3 \\
% w/o tanh gating
%   & 36.2
%   & 28.4
%   & 10.8 \\
% Linear text projector
%   & 71.5
%   & 94.2
%   & 51.5 \\
% % w/o FFN
% %   & 76.3
% %   & 95.9
% %   & 53.4 \\
% \addlinespace
% % \multicolumn{4}{l}{\YA{"Location of cross attn layer"; ADD THESE LINES/subheadings through table}} \\
% \rowcolor{gray!20}
% % \multicolumn{4}{l}{Location of cross-attention layers} \\
% % Early layers (1, 3)
% %   & {XX.X} & {XX.X} & {XX.X} \\
% % Early+mid (1, 3, 5, 7)
% %   & {XX.X} & {XX.X} & {XX.X} \\
% % Mid+late (5, 7, 9 11)
% %   & {XX.X} & {XX.X} & {XX.X} \\
% % Late layers (9, 11)
% %   & {XX.X} & {XX.X} & {XX.X} \\
% % \addlinespace
% % Gaussian point GT
% %   & {XX.X} & {XX.X} & {XX.X} \\
% \bottomrule
% \end{tabular}
% \begin{tablenotes}[flushleft]\footnotesize
% \item Late fusion appends gated cross-attention after the final ViT layer instead of interleaving within layers.
% PODS is evaluated with descriptive text prompts; the frozen baseline uses no text. Both tanh gating on the cross attention blocks and MLP projector RoBERTa are detrimental to model performance.
% % Layer-placement rows vary which ViT blocks receive cross-attention (default: blocks 2, 5, 8, 11).
% % Gaussian point GT replaces the segmentation-mask target with a Gaussian heatmap centered on the bounding-box midpoint ($\sigma{=}1.1$).
% \end{tablenotes}
% \end{threeparttable}
% \end{table}

\begin{wraptable}{r}{0.50\textwidth}
% \begin{table}[t]
\centering
\footnotesize
\tabcolsep=0.5mm
\begin{threeparttable}
\vspace{-9mm}
\caption{%
Architecture ablations on DINOv2 ViT-B/14.
\model{} combines three design choices; each ablation row toggles exactly one (marked \xmark).
}
\vspace{-5mm}
\label{tab:arch_ablation}
\begin{tabular}{@{}ccc
                S[table-format=2.1]
                S[table-format=2.1]
                S[table-format=2.1]@{}}
\toprule
{Early} & {Tanh} & {MLP} & {FG-CLS} & {CORE} & {PODS} \\
{Fusion} & {Gate}  & {Proj.} & $\uparrow$ & $\uparrow$ & $\uparrow$ \\
\midrule
\color{gray} {--}
  & \color{gray} {--}
  & \color{gray} {--}
  & \color{gray} 89.0
  & \color{gray} 43.5
  & \color{gray} 28.6 \\
\rowcolor{\methodcolor}
\cmark & \cmark & \cmark
  & 87.7
  & \textbf{96.0}
  & \textbf{58.1} \\
\addlinespace
\xmark & \cmark & \cmark
  & \textbf{91.8}
  & 93.3
  & 36.6 \\
\cmark & \xmark & \cmark
  & 83.5
  & 94.6
  & 47.1 \\
\cmark & \cmark & \xmark
  & 86.7
  & 95.2
  & 56.4 \\
\bottomrule
\end{tabular}
% \begin{tablenotes}[flushleft]\footnotesize
Top gray row: frozen DINOv2 baseline (no adapter \JR{in Sec. 3, we call the MLP adapter; might be better to use other term here.}).
\xmark\ in Early Fusion means late fusion. %(cross-attention appended after the final ViT layer).
\xmark\ in Tanh Gate means ungated cross-attention.
\xmark\ in MLP Proj.\ means single linear layer instead of two. % replacing the two-layer MLP.
PODS uses descriptive text prompts.
Frozen DINOv2 baseline with no adapter (top row, gray) cannot use text. \JR{merge with first sentence}
% \end{tablenotes}
\end{threeparttable}
\vspace{-6mm}
% \end{table}
\end{wraptable}



We ablate the key architectural and training decisions in \model{} to isolate the contribution of each component.
Unless stated otherwise, all experiments use the DINOv2 ViT-B/14 backbone and the RoBERTa-Large as the text encoder and are trained at $336{\times}336$ resolution with a batch size 12 for 500k iterations (${\sim}$84 H100 GPU-hours) on the full data mixture (\cref{sec:data}), using AdamW with a cosine schedule that warms up to $3{\times}10^{-4}$ over 5k steps, decays to $3{\times}10^{-5}$ by 40k steps, and remains constant thereafter.
% The default gate $\alpha_l = 1$
% \MT{what is the default CA gate value during inference. also, indicate if this is pre/post tanh}\JR{it is the learned one. If we do, then we just vary the learned weights by a certain factor} \MT{not true? then why does fig 1 star appear at $\alpha_l = 1$?}\JR{since a factor of 1.0 is essentially not doing any changes -> default setup. Only if for != 1.0 we change model behavior. }
We evaluate in three domains: fine-grained classification probe accuracy (FG-CLS; \cref{subsec:pareto}), text steerability via CORE retrieval (\cref{subsec:cond_ret_exp}), and personalized object discrimination PR-AUC on PODS with descriptive prompts (\cref{subsec:text_density_exp}). %Additional ablations and scaling experiments are included in \cref{apdx:more_ablations}.


% ------------------------------------------------------------------
\paragraph{Architecture Choices.}
\label{subsec:arch_ablation}


% \cref{tab:arch_ablation} isolates the contribution of each architectural component by modifying one design choice at a time from the full \model{} configuration.
\model{} has three principal design choices:
% \JR{there are many more design choices -> maybe ``principal design choices''; perhaps ``dimensions/components'' is better than choices}:
early fusion of text within the ViT,
tanh-gated CA layers, and
MLP text projector.
\cref{tab:arch_ablation} ablates each component.


%%%%%%%%%%%%%%%%%%%%%%%%%%%%%%%%%%%%%%%%%%%%%%%%%%%%%%%%%%%%%%%%%%%%%%%%%%%%%%%%%%%%%%%%%%%%%%%%%%%%%%%%%%%%%%%%%%
%%%%%%%%%%%%%%%%%%%%%%%%%%%%%%%%%%%%%%%%%%%%%%%%%%%%%%%%%%%%%%%%%%%%%%%%%%%%%%%%%%%%%%%%%%%%%%%%%%%%%%%%%%%%%%%%%%
\textit{Early vs.\ late fusion.}
We interleave gated cross-attention within ViT layers (early fusion).  
Late fusion (\cref{tab:arch_ablation}, row 3) adds text only after the final ViT layer.  
While late fusion (row 3) also achieves high steerability (93.3) and higher FG-CLS (91.8 vs.\ 87.7), it reduces PODS performance dramatically (36.6 vs.\ 58.1).
This illustrates the importance of early fusion for fine-grained vision–language modeling, with the gap vanishing (26.5 vs. 27.9) when conditioning on coarse supercategory prompts.
% With coarse supercategory prompts, the gap vanishes , suggesting that early fusion is important for 

\textit{Role of gating.}
Removing the zero-initialized tanh gate (row 4) reduces FG-CLS, CORE, and PODS by 4.2, 1.4, and 11.0 points below \model{} (row 2), showing that ungated cross-attention disrupts frozen features.
% Each cross-attention block uses a zero-initialized tanh gate (\cref{eq:gated_ca}), starting training from a frozen ViT and gradually incorporating text.  
% Removing gating (row 4) reduces FG-CLS, CORE, and PODS by 4.2, 1.4, and 11.0 points below \model{} (row 2), showing that ungated cross-attention disrupts frozen features.
         
\textit{Language projector.}
A two-layer MLP projects RoBERTa features to the ViT space (\cref{sec:arch}).  
Replacing it with a linear projector (row 5) lowers FG-CLS by 1.0 and PODS by 1.7 points, indicating that the MLP better aligns modalities.

% \paragraph{Early vs.\ late fusion.}
% Our method interleaves gated cross-attention blocks within the ViT layers (early fusion).
% The alternative, late fusion (row 3 of \cref{tab:arch_ablation}), appends them after the final ViT layer adding text information to the ViT output features.
% Late fusion achieves comparable steerability (93.4) and better preserves FG-CLS (86.5 vs.\ 77.0),
% % \MT{what do we train here? only adapter? then this looks pretty good to me! perhaps focus on differentiator that cannot do dense pred}
% but degrades performance on PODS by 10 points (45.3 vs.\ 55.3) when conditioning on descriptive text prompts.
% On PODS with coarse supercategory prompts, however, we find the gap between early and late fusion vanishes, indicating that early fusion is necessary for fine-grained vision-language modeling.


% \paragraph{Role of gating.}
% Each cross-attention block uses a zero-initialized tanh gate (\cref{eq:gated_ca}).
% Training starts as a frozen ViT at initialization and gradually incorporates text.
% Removing gating (row 4) collapses both steerability and FG-CLS by 28.4 and 52.8 points below the frozen baseline respectively. This confirms that unconstrained cross-attention catastrophically interferes with frozen features.

% \paragraph{Language projector.}
% A two-layer MLP projects RoBERTa embeddings into the ViT's feature space prior to cross-attention (\cref{sec:arch}).
% Replacing it with a linear projector (row 5) degrades FG-CLS by 5.5 points and PODS by 3.8 points.
% This illustrates that MLP's additional expressivity aids representation alignment, allowing the cross-attention layers to focus on contextualizing visual features with text rather than bridging the modality gap. \textbf{}
%%%%%%%%%%%%%%%%%%%%%%%%%%%%%%%%%%%%%%%%%%%%%%%%%%%%%%%%%%%%%%%%%%%%%%%%%%%%
%%%%%%%%%%%%%%%%%%%%%%%%%%%%%%%%%%%%%%%%%%%%%%%%%%%%%%%%%%%%%%%%%%%%%%%%%%%%
% while steerability remains comparable (94.2 vs.\ 93.4).


% \MT{this can probably be just moved to supplement}

% TODO: insert results and interpretation once numbers are ready, e.g.:
% Our default spread-out placement achieves the best PODS (XX.X),
% while restricting to late layers preserves FG-CLS but reduces steerability,
% suggesting that text conditioning in early layers is necessary for reshaping
% intermediate representations rather than merely modulating the final output.




% ------------------------------------------------------------------
\paragraph{Generalization Across Pretrained Backbones.}
\label{subsec:cross_backbone}

% \begin{table}[t]
\begin{wraptable}{r}{0.4\textwidth}
\centering
\footnotesize
\tabcolsep=1.2mm
\vspace{-12mm}
\caption{%
\textbf{Steering different ViT pretraining approaches.}
% Results presented as \textbf{base model / late fusion / \model{}}.
\model{} consistently outperforms late fusion, with the largest gains on weaker backbones.
% MVTech PRO measures zero-shot anomaly segmentation, a domain absent from training. 
}
\label{tab:cross_backbone}
\begin{tabular}{l cc >{\columncolor{\methodcolor}}c}
\toprule
\multirow{2}{*}{Backbone} & \multicolumn{3}{c}{CORE $\uparrow$} \\
& Base & Late & \model{} \\
\midrule
DINOv2 & 43.5 & 93.3 & \textbf{96.0} \\
SigLIP & 38.3 & 75.4 & \textbf{91.3} \\
MAE    & 21.8 & 41.0 & \textbf{74.9} \\
\bottomrule
\end{tabular}
\vspace{-6mm}
\end{wraptable}

To validate that our approach generalizes beyond DINOv2, we apply \model{} to two additional ViTs: SigLIP~\cite{zhai2023_siglip}
% (contrastive vision-language pretraining)
and MAE~\cite{he2022mae}.
% (masked autoencoding).
\cref{tab:cross_backbone} compares vanilla, late fusion, and early fusion (\model{}) variants. %\JR{consider adding full ablation table}

Overall, DINOv2 produces the strongest steerable representations, consistent with its richer self-supervised features. SigLIP and MAE similarly benefit from text conditioning but start from weaker baselines.
Early fusion consistently outperforms late fusion across all three ViT families.
The advantage is more pronounced for MAE and SigLIP, where early fusion boosts steerability by 33.9 and 15.9 points, respectively. 
These larger early-fusion gains compared to DINOv2 (2.7) align with our later layer-wise analysis (cf.~\cref{fig:feature_divergence}), where \model{} representations diverge earlier from their base ViTs for MAE and SigLIP backbones, suggesting that language injection is especially beneficial when the underlying visual features are less semantically mature.
\newline \newline
Additional ablations for backbones, model scaling, proxy-task are in \cref{apdx:more_ablations}.